\documentclass{ximera}

\input{../../../preamble.tex}


\definecolor{fvdnavy}{HTML}{071D43}
\definecolor{fvdcyan}{HTML}{20BFC6}
\definecolor{fvdcyanlight}{HTML}{EFFCFB}
\definecolor{fvdmagenta}{HTML}{F10D69}
\definecolor{fvdmagentalight}{HTML}{FFF1F7}
\definecolor{fvdtext}{HTML}{163044}





%=================================================
% Pedagogische omgevingen
% PDF: tcolorbox
% HTML: learning-box
%=================================================

\newenvironment{voorkennis}
{%
    \ifdefined\HCode
        \HCode{%
<section class="learning-box learning-box--prior">
<div class="learning-box__header">
<span class="learning-box__icon" aria-hidden="true"></span>
<span class="learning-box__title">Voorkennis</span>
</div>
<div class="learning-box__content">
        }%
        \begin{itemize}
    \else
       \begin{tcolorbox}[
    enhanced,
    breakable,
    colback=fvdcyanlight,
    colframe=fvdcyan,
    coltext=fvdtext,
    boxrule=0.5pt,
    leftrule=4pt,
    arc=4mm,
    outer arc=4mm,
    left=7mm,
    right=6mm,
    top=5mm,
    bottom=5mm,
    before skip=6mm,
    after skip=6mm,
    title=Voorkennis,
    coltitle=fvdcyan,
    fonttitle=\bfseries\Large,
    attach title to upper={\par\smallskip},
    boxed title style={
        empty,
        left=0mm,
        right=0mm,
        top=0mm,
        bottom=0mm
    },
    overlay={
        \draw[
            fvdcyan,
            line width=0.5pt
        ]
        ([xshift=42mm,yshift=-13mm]frame.north west)
        --
        ([xshift=-7mm,yshift=-13mm]frame.north east);
    }
]
        \begin{itemize}
    \fi
}
{%
    \end{itemize}
    \ifdefined\HCode
        \HCode{%
</div>
</section>
        }%
    \else
        \end{tcolorbox}
    \fi
}


\newenvironment{leerdoelen}
{%
    \ifdefined\HCode
        \HCode{%
<section class="learning-box learning-box--goals">
<div class="learning-box__header">
<span class="learning-box__icon" aria-hidden="true"></span>
<span class="learning-box__title">Leerdoelen</span>
</div>
<div class="learning-box__content">
        }%
        \begin{itemize}
    \else
\begin{tcolorbox}[
    enhanced,
    breakable,
    colback=fvdmagentalight,
    colframe=fvdmagenta,
    coltext=fvdtext,
    boxrule=0.5pt,
    leftrule=4pt,
    arc=4mm,
    outer arc=4mm,
    left=7mm,
    right=6mm,
    top=5mm,
    bottom=5mm,
    before skip=6mm,
    after skip=6mm,
    title=Leerdoelen,
    coltitle=fvdmagenta,
    fonttitle=\bfseries\Large,
    attach title to upper={\par\smallskip},
    boxed title style={
        empty,
        left=0mm,
        right=0mm,
        top=0mm,
        bottom=0mm
    },
    overlay={
        \draw[
            fvdmagenta,
            line width=0.5pt
        ]
        ([xshift=42mm,yshift=-13mm]frame.north west)
        --
        ([xshift=-7mm,yshift=-13mm]frame.north east);
    }
]
        \begin{itemize}
    \fi
}
{%
    \end{itemize}
    \ifdefined\HCode
        \HCode{%
</div>
</section>
        }%
    \else
        \end{tcolorbox}
    \fi
}


\addPrintStyle{../../..}

\title{Oefeningen — Stelsels herbekeken}
\author{Ingmar Herreman}

\begin{document}

% FVD_CHAPTER_DATA_START
% Automatisch gegenereerd uit index.tex.
% Niet handmatig aanpassen.
\fvdchapterdata{1}{Van algebra naar ruimte}{1}
% FVD_CHAPTER_DATA_END







\maketitle

\FVDbeginExercises

% ============================================================
% INLEIDING
% ============================================================

\begin{opmerking}
Dit hoofdstuk is een herhaling van leerstof uit het vijfde jaar.

De bedoeling is niet dat je opnieuw tientallen stelsels volledig
met de hand reduceert. Voor omvangrijke berekeningen mag je ICT gebruiken.

Je moet vooral kunnen:

\begin{itemize}
    \item een stelsel correct invoeren;
    \item de gereduceerde matrix interpreteren;
    \item de rang bepalen;
    \item de oplosbaarheid analyseren;
    \item vrije onbekenden herkennen;
    \item een oplossingsverzameling correct schrijven;
    \item je besluit wiskundig motiveren.
\end{itemize}
\end{opmerking}


% ============================================================
% BASIS
% ============================================================

\FVDsection{Basis — techniek opnieuw activeren}


% ------------------------------------------------------------
% OEFENING 1
% ------------------------------------------------------------

\begin{oefening}[Van stelsel naar matrix]{\oefB\ESS}

Schrijf van elk stelsel de coëfficiëntenmatrix \(A\),
de kolommatrix \(B\) en de uitgebreide matrix \(A_b\).

\begin{question}

\[
\begin{cases}
2x-3y+z=7,\\
x+4y-2z=-1,\\
3x-z=5.
\end{cases}
\]

\begin{oplossing}

De volgorde van de onbekenden is \(x,y,z\).

\[
A=
\begin{bmatrix}
2&-3&1\\
1&4&-2\\
3&0&-1
\end{bmatrix},
\qquad
B=
\begin{bmatrix}
7\\-1\\5
\end{bmatrix}.
\]

Dus

\[
A_b=
\left[
\begin{array}{ccc|c}
2&-3&1&7\\
1&4&-2&-1\\
3&0&-1&5
\end{array}
\right].
\]

Merk op dat in de derde vergelijking de coëfficiënt van \(y\)
gelijk is aan \(0\).

\end{oplossing}

\end{question}


\begin{question}

\[
\begin{cases}
x_1-2x_3+x_4=3,\\
2x_1+x_2+4x_4=-1,\\
x_2+3x_3=5.
\end{cases}
\]

\begin{oplossing}

De volgorde is

\[
x_1,x_2,x_3,x_4.
\]

Daarom

\[
A=
\begin{bmatrix}
1&0&-2&1\\
2&1&0&4\\
0&1&3&0
\end{bmatrix},
\qquad
B=
\begin{bmatrix}
3\\-1\\5
\end{bmatrix}.
\]

De uitgebreide matrix is

\[
A_b=
\left[
\begin{array}{cccc|c}
1&0&-2&1&3\\
2&1&0&4&-1\\
0&1&3&0&5
\end{array}
\right].
\]

\end{oplossing}

\end{question}

\end{oefening}



% ------------------------------------------------------------
% OEFENING 2
% ------------------------------------------------------------

\begin{oefening}[Lees de gereduceerde matrix]{\oefB\ESS}

De volgende matrices zijn reeds met ICT gereduceerd.

Bepaal telkens:

\begin{enumerate}
    \item \(\operatorname{rang}(A)\);
    \item \(\operatorname{rang}(A_b)\);
    \item het aantal vrije onbekenden;
    \item het aantal oplossingen;
    \item indien mogelijk: de oplossingsverzameling.
\end{enumerate}


\begin{question}

\[
\left[
\begin{array}{ccc|c}
1&0&0&2\\
0&1&0&-1\\
0&0&1&5
\end{array}
\right]
\]

\begin{oplossing}

Er zijn drie pivots in \(A\):

\[
\operatorname{rang}(A)=3.
\]

Ook

\[
\operatorname{rang}(A_b)=3.
\]

Er zijn \(n=3\) onbekenden, dus

\[
3-3=0
\]

vrije onbekenden.

Daarom is er precies één oplossing:

\[
V=\{(2,-1,5)\}.
\]

\end{oplossing}

\end{question}


\begin{question}

\[
\left[
\begin{array}{ccc|c}
1&0&2&4\\
0&1&-3&1\\
0&0&0&0
\end{array}
\right]
\]

\begin{oplossing}

\[
\operatorname{rang}(A)
=
\operatorname{rang}(A_b)
=
2.
\]

Omdat \(n=3\), is er

\[
3-2=1
\]

vrije onbekende.

Dus het stelsel heeft oneindig veel oplossingen.

Neem

\[
z=t.
\]

Dan

\[
x+2z=4
\quad\Longrightarrow\quad
x=4-2t
\]

en

\[
y-3z=1
\quad\Longrightarrow\quad
y=1+3t.
\]

Dus

\[
V=
\left\{
(4-2t,1+3t,t)
\mid t\in\mathbb{R}
\right\}.
\]

\end{oplossing}

\end{question}


\begin{question}

\[
\left[
\begin{array}{ccc|c}
1&0&2&4\\
0&1&-3&1\\
0&0&0&1
\end{array}
\right]
\]

\begin{oplossing}

Voor \(A\):

\[
\operatorname{rang}(A)=2.
\]

Voor \(A_b\):

\[
\operatorname{rang}(A_b)=3.
\]

Dus

\[
\operatorname{rang}(A)<\operatorname{rang}(A_b).
\]

De laatste rij betekent bovendien

\[
0=1,
\]

wat onmogelijk is.

Daarom

\[
V=\varnothing.
\]

\end{oplossing}

\end{question}

\end{oefening}



% ------------------------------------------------------------
% OEFENING 3
% ------------------------------------------------------------

\begin{oefening}[Reduceren met ICT en analyseren]{\oefB\ESS}

Gebruik ICT om de uitgebreide matrix van elk stelsel naar gereduceerde
rij-echelonvorm te brengen.

Noteer vervolgens expliciet:

\[
\operatorname{rang}(A),\qquad
\operatorname{rang}(A_b),
\]

het aantal vrije onbekenden en je besluit over de oplosbaarheid.


\begin{question}

\[
\begin{cases}
2x+3y+z=11,\\
3x+y+2z=11,\\
x+2y+3z=14.
\end{cases}
\]

\begin{oplossing}

Met ICT:

\[
A_b
\sim
\left[
\begin{array}{ccc|c}
1&0&0&1\\
0&1&0&2\\
0&0&1&3
\end{array}
\right].
\]

Dus

\[
\operatorname{rang}(A)
=
\operatorname{rang}(A_b)
=
3=n.
\]

Er zijn geen vrije onbekenden.

Het stelsel heeft precies één oplossing:

\[
V=\{(1,2,3)\}.
\]

\end{oplossing}

\end{question}


\begin{question}

\[
\begin{cases}
x-2y=1,\\
2x+y-z=1,\\
3x-y-z=2.
\end{cases}
\]

\begin{oplossing}

Met ICT:

\[
A_b
\sim
\left[
\begin{array}{ccc|c}
1&0&-\frac25&\frac35\\
0&1&-\frac15&-\frac15\\
0&0&0&0
\end{array}
\right].
\]

Daarom

\[
\operatorname{rang}(A)
=
\operatorname{rang}(A_b)
=
2<3.
\]

Er is één vrije onbekende.

Neem

\[
z=t.
\]

Dan

\[
x=\frac35+\frac25t
\]

en

\[
y=-\frac15+\frac15t.
\]

Dus

\[
V=
\left\{
\left(
\frac35+\frac25t,
-\frac15+\frac15t,
t
\right)
\mid t\in\mathbb{R}
\right\}.
\]

\end{oplossing}

\end{question}


\begin{question}

\[
\begin{cases}
2x+2y+4z=13,\\
y+z=5,\\
x+2y+3z=8.
\end{cases}
\]

\begin{oplossing}

Met ICT:

\[
A_b
\sim
\left[
\begin{array}{ccc|c}
1&0&1&0\\
0&1&1&0\\
0&0&0&1
\end{array}
\right].
\]

Voor de coëfficiëntenmatrix geldt

\[
\operatorname{rang}(A)=2,
\]

maar voor de uitgebreide matrix

\[
\operatorname{rang}(A_b)=3.
\]

Dus

\[
\operatorname{rang}(A)
<
\operatorname{rang}(A_b).
\]

Het stelsel is strijdig:

\[
V=\varnothing.
\]

\end{oplossing}

\end{question}

\end{oefening}



% ------------------------------------------------------------
% OEFENING 4
% ------------------------------------------------------------

\begin{oefening}[Vrije onbekenden]{\oefB\ESS}

Voor elk van de volgende gereduceerde matrices:

\begin{enumerate}
    \item duid de pivotkolommen aan;
    \item geef de vrije onbekenden;
    \item bepaal het aantal vrije parameters;
    \item schrijf de volledige oplossingsverzameling.
\end{enumerate}


\begin{question}

\[
\left[
\begin{array}{ccc|c}
1&0&2&3\\
0&1&-1&4\\
0&0&0&0
\end{array}
\right]
\]

\begin{oplossing}

De pivots liggen in kolommen \(1\) en \(2\).

Dus \(x\) en \(y\) zijn pivotvariabelen en \(z\) is vrij.

Neem

\[
z=t.
\]

Dan

\[
x=3-2t,
\qquad
y=4+t.
\]

Dus

\[
V=
\left\{
(3-2t,4+t,t)
\mid t\in\mathbb{R}
\right\}.
\]

Equivalent:

\[
\begin{pmatrix}
x\\y\\z
\end{pmatrix}
=
\begin{pmatrix}
3\\4\\0
\end{pmatrix}
+
t
\begin{pmatrix}
-2\\1\\1
\end{pmatrix}.
\]

\end{oplossing}

\end{question}


\begin{question}

\[
\left[
\begin{array}{cccc|c}
1&0&-2&3&1\\
0&1&1&-1&2\\
0&0&0&0&0
\end{array}
\right]
\]

\begin{oplossing}

De pivots liggen in kolommen \(1\) en \(2\).

Dus

\[
x_3=s,
\qquad
x_4=t
\]

zijn vrij.

Uit de eerste rij:

\[
x_1-2x_3+3x_4=1
\]

volgt

\[
x_1=1+2s-3t.
\]

Uit de tweede rij:

\[
x_2+x_3-x_4=2
\]

volgt

\[
x_2=2-s+t.
\]

Dus

\[
V=
\left\{
(1+2s-3t,\,2-s+t,\,s,\,t)
\mid s,t\in\mathbb{R}
\right\}.
\]

\end{oplossing}

\end{question}

\end{oefening}



% ------------------------------------------------------------
% OEFENING 5
% ------------------------------------------------------------

\begin{oefening}[Homogene stelsels]{\oefB\ESS}

Gebruik ICT waar nodig.


\begin{question}

Beschouw

\[
\begin{cases}
2x+y+3z=0,\\
y+2z=0,\\
-x+5y=0.
\end{cases}
\]

Heeft dit stelsel niet-triviale oplossingen?

Motiveer met behulp van de rang.

\begin{oplossing}

De reductie geeft

\[
A_b
\sim
\left[
\begin{array}{ccc|c}
1&0&0&0\\
0&1&0&0\\
0&0&1&0
\end{array}
\right].
\]

Dus

\[
\operatorname{rang}(A)=3=n.
\]

Er zijn geen vrije onbekenden.

Het stelsel heeft enkel de triviale oplossing:

\[
V=\{(0,0,0)\}.
\]

\end{oplossing}

\end{question}


\begin{question}

Beschouw

\[
\begin{cases}
x+3y-z=0,\\
2x+6y-3z=0,\\
x+3y+3z=0.
\end{cases}
\]

Bepaal de oplossingsverzameling.

\begin{oplossing}

Met ICT:

\[
A_b
\sim
\left[
\begin{array}{ccc|c}
1&3&0&0\\
0&0&1&0\\
0&0&0&0
\end{array}
\right].
\]

Dus

\[
\operatorname{rang}(A)=2<3.
\]

Er is één vrije onbekende.

Uit de tweede rij:

\[
z=0.
\]

Neem

\[
y=t.
\]

Dan

\[
x=-3t.
\]

Dus

\[
V=
\left\{
(-3t,t,0)
\mid t\in\mathbb{R}
\right\}.
\]

\end{oplossing}

\end{question}

\end{oefening}



% ============================================================
% VERDIEPING
% ============================================================

\FVDsection{Verdieping — analyseren en redeneren}


% ------------------------------------------------------------
% OEFENING 6
% ------------------------------------------------------------

\begin{oefening}[Beslis zonder verder te rekenen]{\oefU\ESS}

Voor elk geval is de reductie voldoende ver uitgevoerd om de oplosbaarheid
te kunnen bepalen.

Je mag niet verder reduceren.

Bepaal het aantal oplossingen en motiveer.


\begin{question}

\[
\left[
\begin{array}{ccc|c}
1&2&-1&4\\
0&1&3&2\\
0&0&5&-10
\end{array}
\right].
\]

\begin{oplossing}

Alle drie de coëfficiëntenkolommen bevatten een pivot.

Daarom

\[
\operatorname{rang}(A)
=
\operatorname{rang}(A_b)
=
3=n.
\]

Het stelsel heeft precies één oplossing.

Het is niet nodig de oplossing zelf uit te rekenen.

\end{oplossing}

\end{question}


\begin{question}

\[
\left[
\begin{array}{ccc|c}
1&2&-1&4\\
0&1&3&2\\
0&0&0&-6
\end{array}
\right].
\]

\begin{oplossing}

De laatste rij geeft

\[
0=-6.
\]

Dat is onmogelijk.

Equivalent:

\[
\operatorname{rang}(A)=2
\]

maar

\[
\operatorname{rang}(A_b)=3.
\]

Dus het stelsel heeft geen oplossing.

\end{oplossing}

\end{question}


\begin{question}

\[
\left[
\begin{array}{cccc|c}
1&2&0&-1&5\\
0&0&1&3&2\\
0&0&0&0&0
\end{array}
\right].
\]

\begin{oplossing}

Er zijn twee pivots:

\[
\operatorname{rang}(A)
=
\operatorname{rang}(A_b)
=
2.
\]

Er zijn vier onbekenden.

Dus

\[
4-2=2
\]

vrije onbekenden.

Het stelsel heeft oneindig veel oplossingen met twee vrije parameters.

\end{oplossing}

\end{question}

\end{oefening}



% ------------------------------------------------------------
% OEFENING 7
% ------------------------------------------------------------

\begin{oefening}[Vind de fout in de redenering]{\oefU\ESS}

Een leerling krijgt na reductie

\[
\left[
\begin{array}{ccc|c}
1&0&1&-2\\
0&1&1&5\\
0&0&0&7
\end{array}
\right].
\]

Hij schrijft:

\[
\operatorname{rang}(A)=3,
\qquad
\operatorname{rang}(A_b)=2.
\]

Daarom heeft het stelsel geen oplossing.

Analyseer deze redenering.

\begin{oplossing}

De conclusie ``geen oplossing'' is juist, maar de rangen zijn fout.

Voor de coëfficiëntenmatrix

\[
A=
\begin{bmatrix}
1&0&1\\
0&1&1\\
0&0&0
\end{bmatrix}
\]

zijn er twee pivots. Dus

\[
\operatorname{rang}(A)=2.
\]

In de uitgebreide matrix ontstaat door de laatste kolom een derde pivot:

\[
\operatorname{rang}(A_b)=3.
\]

Dus

\[
\operatorname{rang}(A)
<
\operatorname{rang}(A_b),
\]

en daarom is het stelsel strijdig.

Bovendien geldt algemeen altijd

\[
\operatorname{rang}(A_b)
\geq
\operatorname{rang}(A),
\]

omdat \(A_b\) ontstaat door aan \(A\) een kolom toe te voegen.

\end{oplossing}

\end{oefening}



% ------------------------------------------------------------
% OEFENING 8
% ------------------------------------------------------------

\begin{oefening}[Zelfde coëfficiënten, ander rechterlid]{\oefU\ESS}

Beschouw de twee gereduceerde uitgebreide matrices

\[
M_1=
\left[
\begin{array}{ccc|c}
1&0&2&3\\
0&1&-1&4\\
0&0&0&0
\end{array}
\right]
\]

en

\[
M_2=
\left[
\begin{array}{ccc|c}
1&0&2&0\\
0&1&-1&0\\
0&0&0&1
\end{array}
\right].
\]

\begin{question}

De coëfficiëntenmatrices hebben dezelfde rang.
Toch is de oplosbaarheid volledig verschillend.

Verklaar.

\begin{oplossing}

Voor beide matrices is

\[
\operatorname{rang}(A)=2.
\]

Bij \(M_1\):

\[
\operatorname{rang}(A_b)=2,
\]

dus

\[
\operatorname{rang}(A)
=
\operatorname{rang}(A_b)
<
3.
\]

Het stelsel heeft oneindig veel oplossingen.

Bij \(M_2\) geeft de laatste rij

\[
0=1.
\]

Daarom

\[
\operatorname{rang}(A_b)=3.
\]

Dus

\[
\operatorname{rang}(A)
<
\operatorname{rang}(A_b),
\]

en het stelsel heeft geen oplossing.

De rang van \(A\) alleen is dus onvoldoende om de oplosbaarheid
van een niet-homogeen stelsel te bepalen.

\end{oplossing}

\end{question}

\end{oefening}



% ------------------------------------------------------------
% OEFENING 9
% ------------------------------------------------------------

\begin{oefening}[Hoeveel vrijheid blijft over?]{\oefU\ESS}

Een oplosbaar lineair stelsel bevat \(7\) onbekenden.

De coëfficiëntenmatrix heeft rang \(4\).

\begin{question}

Hoeveel vrije parameters heeft de oplossingsverzameling?

\[
\answer{3}
\]

\begin{oplossing}

Omdat het stelsel oplosbaar is,

\[
\operatorname{rang}(A)
=
\operatorname{rang}(A_b).
\]

Het aantal vrije parameters is

\[
n-\operatorname{rang}(A)
=
7-4
=
3.
\]

\end{oplossing}

\end{question}


\begin{question}

Heeft het stelsel één of oneindig veel oplossingen?

\begin{multipleChoice}
    \choice{geen oplossing}
    \choice{precies één oplossing}
    \choice[correct]{oneindig veel oplossingen}
\end{multipleChoice}

\begin{oplossing}

Er zijn drie vrije parameters. Elke parameter kan vrij variëren in
\(\mathbb{R}\), dus er zijn oneindig veel oplossingen.

\end{oplossing}

\end{question}

\end{oefening}



% ------------------------------------------------------------
% OEFENING 10
% ------------------------------------------------------------

\begin{oefening}[Van parameter naar ruimte]{\oefU\ESS}

Beschouw

\[
V=
\left\{
(3-2t,4+t,t)
\mid t\in\mathbb{R}
\right\}.
\]

\begin{question}

Schrijf deze oplossingsverzameling in de vorm

\[
\begin{pmatrix}
x\\y\\z
\end{pmatrix}
=
\begin{pmatrix}
x_0\\y_0\\z_0
\end{pmatrix}
+
t
\begin{pmatrix}
a\\b\\c
\end{pmatrix}.
\]

\begin{oplossing}

\[
\begin{pmatrix}
x\\y\\z
\end{pmatrix}
=
\begin{pmatrix}
3\\4\\0
\end{pmatrix}
+
t
\begin{pmatrix}
-2\\1\\1
\end{pmatrix}.
\]

\end{oplossing}

\end{question}


\begin{question}

Bereken de punten van de oplossingsverzameling voor

\[
t=0,\qquad t=1,\qquad t=-1.
\]

\begin{oplossing}

Voor \(t=0\):

\[
(3,4,0).
\]

Voor \(t=1\):

\[
(1,5,1).
\]

Voor \(t=-1\):

\[
(5,3,-1).
\]

\end{oplossing}

\end{question}


\begin{question}

Wat voor meetkundige figuur vormen alle oplossingen?

\begin{multipleChoice}
    \choice{één punt}
    \choice[correct]{een rechte}
    \choice{een vlak}
    \choice{de volledige ruimte}
\end{multipleChoice}

\begin{oplossing}

Er is één vrije parameter.

Wanneer \(t\) varieert, vertrekken we vanuit het vaste punt

\[
(3,4,0)
\]

en bewegen we in de vaste richting

\[
(-2,1,1).
\]

Daarom vormt de oplossingsverzameling een rechte.

\end{oplossing}

\end{question}

\end{oefening}



% ------------------------------------------------------------
% OEFENING 11
% ------------------------------------------------------------

\begin{oefening}[Rang en meetkundige betekenis]{\oefU\ESS}

Een oplosbaar stelsel heeft drie onbekenden \(x,y,z\).

Vul de tabel aan.

\[
\begin{array}{c|c|c|c}
\operatorname{rang}(A)
&
\text{vrije parameters}
&
\text{aantal oplossingen}
&
\text{meetkundige vorm}
\\
\hline
3 & ? & ? & ?\\
2 & ? & ? & ?\\
1 & ? & ? & ?\\
0 & ? & ? & ?
\end{array}
\]

\begin{oplossing}

\[
\begin{array}{c|c|c|c}
\operatorname{rang}(A)
&
\text{vrije parameters}
&
\text{aantal oplossingen}
&
\text{meetkundige vorm}
\\
\hline
3&0&1&\text{punt}\\
2&1&\infty&\text{rechte}\\
1&2&\infty&\text{vlak}\\
0&3&\infty&\mathbb{R}^{3}
\end{array}
\]

De tabel geldt alleen wanneer het stelsel oplosbaar is.

\end{oplossing}

\end{oefening}



% ------------------------------------------------------------
% OEFENING 12
% ------------------------------------------------------------

\begin{oefening}[Een homogeen stelsel analyseren zonder het op te lossen]{\oefU}

Een homogeen stelsel heeft \(5\) onbekenden en

\[
\operatorname{rang}(A)=3.
\]

Beantwoord zonder het stelsel expliciet op te lossen.

\begin{question}

Kan het stelsel strijdig zijn?

\begin{multipleChoice}
    \choice{ja}
    \choice[correct]{nee}
\end{multipleChoice}

\begin{oplossing}

Nee.

Een homogeen stelsel heeft altijd minstens de nuloplossing.

\end{oplossing}

\end{question}


\begin{question}

Hoeveel vrije parameters zijn er?

\[
\answer{2}
\]

\begin{oplossing}

\[
5-3=2.
\]

\end{oplossing}

\end{question}


\begin{question}

Bestaan er niet-triviale oplossingen?

\begin{multipleChoice}
    \choice[correct]{ja}
    \choice{nee}
\end{multipleChoice}

\begin{oplossing}

Ja.

Omdat

\[
\operatorname{rang}(A)=3<5,
\]

zijn er twee vrije parameters. De nuloplossing is dus niet de enige
oplossing.

\end{oplossing}

\end{question}

\end{oefening}



% ------------------------------------------------------------
% OEFENING 13
% ------------------------------------------------------------

\begin{oefening}[Kies de juiste informatie]{\oefU}

Voor een stelsel met vier onbekenden is gegeven:

\[
\operatorname{rang}(A)=3.
\]

Een leerling beweert:

\begin{center}
``Het stelsel heeft dus oneindig veel oplossingen.''
\end{center}

Is dit besluit correct?

\begin{oplossing}

Niet noodzakelijk.

Uit

\[
\operatorname{rang}(A)=3<4
\]

weten we alleen dat er, \emph{als het stelsel oplosbaar is},
één vrije parameter zal zijn.

We moeten eerst ook

\[
\operatorname{rang}(A_b)
\]

kennen.

Er zijn twee mogelijkheden:

\[
\operatorname{rang}(A_b)=3
\]

geeft oneindig veel oplossingen, maar

\[
\operatorname{rang}(A_b)=4
\]

geeft geen oplossing.

De gegeven informatie is dus onvoldoende.

\end{oplossing}

\end{oefening}



% ============================================================
% GEVORDERD
% ============================================================

\FVDsection{Gevorderd — combineren en generaliseren}


% ------------------------------------------------------------
% OEFENING 14
% ------------------------------------------------------------

\begin{oefening}[Een parameter in het rechterlid]{\oefG}

Beschouw

\[
\begin{cases}
x+y+z=2,\\
2x+2y+2z=4,\\
x+y+z=k.
\end{cases}
\]

Onderzoek afhankelijk van \(k\) de oplosbaarheid van het stelsel.

\begin{oplossing}

De tweede vergelijking is tweemaal de eerste.

De derde vergelijking heeft hetzelfde linkerlid als de eerste.

Als

\[
k=2,
\]

zijn de eerste en derde vergelijking identiek en geldt

\[
\operatorname{rang}(A)
=
\operatorname{rang}(A_b)
=
1<3.
\]

Er zijn dan twee vrije parameters en dus oneindig veel oplossingen.

Als

\[
k\neq2,
\]

vragen de eerste en derde vergelijking tegelijk dat

\[
x+y+z=2
\]

en

\[
x+y+z=k.
\]

Dat is onmogelijk.

Na reductie ontstaat een rij van de vorm

\[
[0\ 0\ 0\mid k-2].
\]

Dus voor

\[
k\neq2
\]

geldt

\[
\operatorname{rang}(A)
<
\operatorname{rang}(A_b),
\]

en heeft het stelsel geen oplossing.

Samengevat:

\[
\boxed{
\begin{array}{ll}
k=2 &: \text{oneindig veel oplossingen},\\[1mm]
k\neq2 &: \text{geen oplossing}.
\end{array}
}
\]

\end{oplossing}

\end{oefening}



% ------------------------------------------------------------
% OEFENING 15
% ------------------------------------------------------------

\begin{oefening}[Kan dit bestaan?]{\oefG}

Bepaal telkens of de beschreven situatie mogelijk is.
Motiveer.


\begin{question}

Een stelsel met drie onbekenden waarvoor

\[
\operatorname{rang}(A)=3
\]

en

\[
\operatorname{rang}(A_b)=2.
\]

\begin{oplossing}

Dit is onmogelijk.

Omdat \(A_b\) ontstaat door aan \(A\) één kolom toe te voegen, geldt altijd

\[
\operatorname{rang}(A_b)
\geq
\operatorname{rang}(A).
\]

\end{oplossing}

\end{question}


\begin{question}

Een stelsel met vier onbekenden waarvoor

\[
\operatorname{rang}(A)
=
\operatorname{rang}(A_b)
=
3
\]

en precies één oplossing.

\begin{oplossing}

Dit is onmogelijk.

Er zijn vier onbekenden maar slechts drie pivots.

Dus

\[
4-3=1
\]

onbekende is vrij.

Het stelsel heeft daarom oneindig veel oplossingen.

\end{oplossing}

\end{question}


\begin{question}

Een homogeen stelsel met zes onbekenden, rang \(6\)
en alleen de nuloplossing.

\begin{oplossing}

Dit is mogelijk.

Omdat

\[
\operatorname{rang}(A)=6=n,
\]

zijn alle onbekenden pivotvariabelen en zijn er geen vrije parameters.

Het homogene stelsel heeft dan enkel

\[
X=0.
\]

\end{oplossing}

\end{question}

\end{oefening}



% ------------------------------------------------------------
% OEFENING 16
% ------------------------------------------------------------

\begin{oefening}[Van een rechte naar een stelsel]{\oefG}

Gegeven is de verzameling

\[
V=
\left\{
(1+2t,-1+t,3-t)
\mid t\in\mathbb{R}
\right\}.
\]

\begin{question}

Toon aan dat voor elk punt van \(V\)

\[
x-2y=3
\]

geldt.

\begin{oplossing}

Voor een punt van \(V\) geldt

\[
x=1+2t,
\qquad
y=-1+t.
\]

Dus

\[
x-2y
=
1+2t-2(-1+t)
=
1+2t+2-2t
=
3.
\]

\end{oplossing}

\end{question}


\begin{question}

Zoek nog een lineaire vergelijking in \(x,y,z\)
die door elk punt van \(V\) wordt voldaan.

\begin{oplossing}

Er zijn meerdere mogelijkheden.

We hebben

\[
y=-1+t
\]

en

\[
z=3-t.
\]

Daarom

\[
y+z
=
(-1+t)+(3-t)
=
2.
\]

Een mogelijke tweede vergelijking is dus

\[
y+z=2.
\]

\end{oplossing}

\end{question}


\begin{question}

Beschouw nu het stelsel

\[
\begin{cases}
x-2y=3,\\
y+z=2.
\end{cases}
\]

Toon met behulp van rang en vrije parameters aan dat het
oneindig veel oplossingen heeft.

\begin{oplossing}

De uitgebreide matrix is

\[
\left[
\begin{array}{ccc|c}
1&-2&0&3\\
0&1&1&2
\end{array}
\right].
\]

Er zijn twee pivots:

\[
\operatorname{rang}(A)
=
\operatorname{rang}(A_b)
=
2.
\]

Er zijn drie onbekenden, dus

\[
3-2=1
\]

vrije parameter.

Het stelsel heeft daarom oneindig veel oplossingen.

Neem

\[
z=s.
\]

Dan

\[
y=2-s
\]

en

\[
x=3+2y
=7-2s.
\]

Dus

\[
V=
\left\{
(7-2s,2-s,s)
\mid s\in\mathbb{R}
\right\}.
\]

Dit is dezelfde rechte als de oorspronkelijke verzameling,
maar met een andere parametrisatie.

Neem namelijk

\[
s=3-t.
\]

Dan

\[
(7-2s,2-s,s)
=
(1+2t,-1+t,3-t).
\]

\end{oplossing}

\end{question}

\end{oefening}



% ------------------------------------------------------------
% OEFENING 17
% ------------------------------------------------------------

\begin{oefening}[Redeneren over rang]{\oefG}

Leg uit waarom voor elke matrix \(A\) en haar uitgebreide matrix \(A_b\)

\[
\operatorname{rang}(A_b)
\geq
\operatorname{rang}(A)
\]

moet gelden.

\begin{oplossing}

De uitgebreide matrix \(A_b\) ontstaat door aan \(A\) één extra kolom
toe te voegen.

Alle kolommen die reeds in \(A\) aanwezig waren, blijven dus aanwezig
in \(A_b\).

Door een kolom toe te voegen kunnen bestaande onafhankelijke
kolommen niet plots afhankelijk worden.

De maximale hoeveelheid lineair onafhankelijke informatie kan daarom
niet afnemen.

Dus

\[
\operatorname{rang}(A_b)
\geq
\operatorname{rang}(A).
\]

De extra kolom kan de rang:

\begin{itemize}
    \item onveranderd laten;
    \item of met één verhogen.
\end{itemize}

Daarom is ook

\[
\operatorname{rang}(A_b)
\leq
\operatorname{rang}(A)+1.
\]

\end{oplossing}

\end{oefening}



% ============================================================
% SLOTOPDRACHT
% ============================================================

\FVDsection{Eindcheck}


\begin{oefening}[Alles samen]{\oefU\ESS}

Een stelsel met onbekenden \(x,y,z\) wordt met ICT gereduceerd tot

\[
\left[
\begin{array}{ccc|c}
1&0&-2&5\\
0&1&3&-1\\
0&0&0&0
\end{array}
\right].
\]

Beantwoord zonder bijkomende ICT.


\begin{question}

Bepaal

\[
\operatorname{rang}(A)
\qquad\text{en}\qquad
\operatorname{rang}(A_b).
\]

\begin{oplossing}

\[
\operatorname{rang}(A)
=
\operatorname{rang}(A_b)
=
2.
\]

\end{oplossing}

\end{question}


\begin{question}

Hoeveel vrije onbekenden zijn er?

\[
\answer{1}
\]

\end{question}


\begin{question}

Hoeveel oplossingen heeft het stelsel?

\begin{multipleChoice}
    \choice{geen}
    \choice{precies één}
    \choice[correct]{oneindig veel}
\end{multipleChoice}

\end{question}


\begin{question}

Schrijf de oplossingsverzameling.

\begin{oplossing}

De derde kolom bevat geen pivot, dus neem

\[
z=t.
\]

Dan

\[
x-2z=5
\quad\Longrightarrow\quad
x=5+2t
\]

en

\[
y+3z=-1
\quad\Longrightarrow\quad
y=-1-3t.
\]

Dus

\[
V=
\left\{
(5+2t,-1-3t,t)
\mid t\in\mathbb{R}
\right\}.
\]

\end{oplossing}

\end{question}


\begin{question}

Schrijf de oplossing in vectorvorm.

\begin{oplossing}

\[
\begin{pmatrix}
x\\y\\z
\end{pmatrix}
=
\begin{pmatrix}
5\\-1\\0
\end{pmatrix}
+
t
\begin{pmatrix}
2\\-3\\1
\end{pmatrix}.
\]

\end{oplossing}

\end{question}


\begin{question}

Welke meetkundige figuur beschrijft deze oplossingsverzameling?

\begin{multipleChoice}
    \choice{een punt}
    \choice[correct]{een rechte}
    \choice{een vlak}
    \choice{de volledige ruimte}
\end{multipleChoice}

\begin{oplossing}

Er is één vrije parameter.

De oplossing heeft de vorm

\[
P=P_0+t\vec{v}.
\]

Dat beschrijft een rechte in de ruimte.

Dit is precies de vorm die in een volgend hoofdstuk opnieuw zal verschijnen.

\end{oplossing}

\end{question}

\end{oefening}



% ============================================================
% ESSENTIËLE OEFENINGEN
% ============================================================

\FVDsection{Essentiële oefeningen}

\begin{opmerking}
Om deze opfrissing voldoende te beheersen voor het vervolg van de module,
moet je minstens de volgende oefeningen zelfstandig aankunnen:

\begin{itemize}
    \item \textbf{Oefening 1:} een stelsel correct naar matrixnotatie omzetten;
    \item \textbf{Oefening 2:} rang en oplosbaarheid uit een gereduceerde matrix aflezen;
    \item \textbf{Oefening 3:} een stelsel met ICT reduceren en het resultaat analyseren;
    \item \textbf{Oefening 4:} vrije onbekenden herkennen en parametriseren;
    \item \textbf{Oefening 5:} homogene stelsels analyseren;
    \item \textbf{Oefening 6:} oplosbaarheid bepalen zonder onnodig verder te rekenen;
    \item \textbf{Oefening 7:} fouten in een rangredenering herkennen;
    \item \textbf{Oefening 8:} \(\operatorname{rang}(A)\) en
          \(\operatorname{rang}(A_b)\) correct onderscheiden;
    \item \textbf{Oefening 9:} rang koppelen aan het aantal vrije parameters;
    \item \textbf{Oefening 10:} een geparametriseerde oplossing in vectorvorm schrijven;
    \item \textbf{Oefening 11:} rang koppelen aan de meetkundige dimensie
          van een oplosbare oplossingsverzameling;
    \item \textbf{Eindcheck:} de volledige redenering zelfstandig uitvoeren.
\end{itemize}

De oefeningen uit de categorie \textbf{Gevorderd} zijn niet noodzakelijk
voor het minimumtraject, maar bereiden sterk voor op verdere
ruimtemeetkunde en lineaire algebra.
\end{opmerking}


\end{document}
